\documentclass[a4paper,10pt]{article}
\usepackage[margin=1in]{geometry}
\usepackage{enumitem}
\usepackage{titlesec}
\usepackage[utf8]{inputenc}
\usepackage{hyperref}

\titleformat{\section}{\large\bfseries}{}{0em}{}[\titlerule]
\setlength{\parindent}{0pt}

\begin{document}

\begin{center}
    \textbf{\Large Rody Vilchez}\\
    \vspace{0.3em}
    \normalsize
    Estudiante de Ciencias de la Computación | Data \& AI Enthusiast — Perú\\
    Email: {rody.vilchez00@gmail.com} \quad
    Tel: +51 987 082 126\\
    LinkedIn: \href{https://www.linkedin.com/in/r0sewt/}{Rody Vilchez} \quad
    WebPage: \href{https://rosewt.dev}{rosewt.dev}
    GitHub: \href{https://github.com/r0sewt}{R0SEWT}


\end{center}

\vspace{1em}

\section*{Educación}
\textbf{Ciencias de la Computación (en curso)} \hfill \textit{UPC, Perú}\\
8vo ciclo — Décimo superior\\
Graduación estimada: 2026-2


\section*{Proyectos}

\textbf{Clasificador de Emociones con NAO} \hfill \href{https://github.com/R0SEWT/nao-cnn-emotion}{[GitHub]}\\
Sistema de detección de emociones para robot NAO usando Redes Neuronales Convolucionales (CNNs) y keypoints vía API Flask. Diseñado para interacción expresiva con niños con TEA.\\
\textit{Herramientas: Python, PyTorch, Flask, TensorFlow, MediaPipe, Optuna}

\vspace{0.7em}

\textbf{Análisis de Precios de Viviendas en Dinamarca} \hfill \href{https://github.com/R0SEWT/Denmark-HousePrices-Analysis}{[GitHub]}\\
Modelado exploratorio y predictivo de precios residenciales (1992–2024) con pipelines de Big Data. Uso de XGBoost + Optuna y H2O AutoML, ingeniería de características e interpretabilidad con SHAP.\\
\textit{Herramientas: Python, H2O.ai, SHAP, XGBoost, Optuna, Pandas}

\vspace{0.7em}

\textbf{SpotiOrganizer} \hfill \href{https://github.com/R0SEWT/spotify-playlist-organizer}{[GitHub]}\\
Aplicación que organiza canciones según tempo y estado de ánimo, generando playlists personalizadas automáticamente. Prototipo de sistema de recomendación musical.\\
\textit{Herramientas: Python, Spotify API, scikit-learn, clustering}

\section*{Habilidades}
\begin{itemize}[leftmargin=1.5em, itemsep=0.2em]
    \item \textbf{Lenguajes \& Querying:} Python, SQL, R, C++
    \item \textbf{Data \& ML:} Pandas, scikit-learn, H2O.ai, SHAP, XGBoost, Optuna, Matplotlib, Seaborn, Jupyter
    \item \textbf{Herramientas:} Git, Docker, Excel , Power BI, Google Cloud, Azure (básico)
    \item \textbf{Otros:} FastAPI, Flask, APIs REST, Automatización de procesos con Python
\end{itemize}

\section*{Experiencia}
\textbf{QA Trainee (Practicante)} \hfill \textit{Dic 2024 – Presente}\\
\textit{Visma Latam}
\begin{itemize}[leftmargin=1.5em]
    \item Automatización de validaciones y pruebas de interfaz con Cypress, integrando extracción y verificación de datos para asegurar calidad y consistencia.
    \item Diseño, ejecución y documentación de casos de prueba funcionales, incluyendo validaciones de datos y reportes.
    \item Registro y seguimiento de bugs en Jira, colaborando con desarrollo para resolver incidencias basadas en métricas.
    \item Creación de scripts para control de calidad previo a importación masiva de casos de prueba.
\end{itemize}

\section*{Certificaciones}
\begin{itemize}[leftmargin=1.5em]

    \item \textbf{AI Engineer for Data Scientists Associate certificate} - DataCamp 
    \item \textbf{Especialización en Machine Learning en Google Cloud} — Google Cloud Skills Boost
    \item \textbf{Certificación Profesional en Análisis de Datos} — Google
    \item \textbf{Python for Everybody} — University of Michigan
    \item \textbf{Winter Camp AI 2025} — PUCP
    \item \textbf{Fundamentos de IA Generativa} — IBM
    \item \textbf{Seguridad informática: defensa contra las artes oscuras digitales} — Google
    \item \textbf{IA centrada en el ser humano} — Tecnológico de Monterrey
    \item \textit{Verificación:} \href{https://www.coursera.org/user/5c475ee3e93cd0579b7a95bb0a6deaf8}{coursera.org/user/5c475ee3e93cd0579b7a95bb0a6deaf8}
\end{itemize}


\section*{Voluntariado}
\textbf{Asociación KP — Voluntario} \hfill \textit{2022–2023}\\
\begin{itemize}[leftmargin=1.5em]
    \item 95 horas de voluntariado en bienestar animal: apoyo a refugios, campañas de esterilización y educación comunitaria en Lima.
\end{itemize}

\end{document}
